% Libro 1 · Mundo 2 · Aritmética y promedios (app: pu-aritmetica)
\mundo{Aritmética y promedios}{pu-aritmetica}{Aquí los números ya no se operan sueltos: viven dentro de una historia.
Alguien reparte, gasta, promedia o esconde cifras, y hay que reconstruir lo que pasó. El secreto de este mundo es
\emph{trabajar hacia atrás} y \emph{representar} con barras y tablas antes de escribir una sola ecuación.}{%
\item \textbf{2.1} Fracciones de lo que queda: trabajar hacia atrás
\item \textbf{2.2} Promedios: el nivelador
\item \textbf{2.3} Criptoaritmética: cuentas con letras
\item \textbf{2.4} Máximos y mínimos: el peor y el mejor caso}

% ------------------------------------------------------------------
\leccion{pu-ar-1}{Fracciones de lo que queda}{Juegos de la barra que se parte y del camino de regreso.}

\begin{reto}
Un vendedor de mangos vende en la mañana la mitad de sus mangos más 2; en la tarde, la mitad de los que le quedan más 2,
y se queda con 5 mangos. ¿Con cuántos empezó?
\opciones{28}{32}{36}{40}
\end{reto}

\subsection{La idea: deshacer el camino}
Cuando una historia termina con un dato conocido («le quedan 5») y cada paso es una operación sencilla, lo más rápido
es caminar \clave{de atrás hacia adelante}, haciendo en cada paso la operación \emph{contraria}: lo que se restó se suma,
lo que se partió a la mitad se duplica.

\begin{center}
\begin{tikzpicture}[x=1mm,y=1mm,font=\small]
  \node[draw=zolive,rounded corners=2pt,minimum width=14mm,minimum height=7mm] (a) at (0,0) {$x$};
  \node[draw=zolive,rounded corners=2pt,minimum width=14mm,minimum height=7mm] (b) at (32,0) {$\tfrac{x}{2} - 2$};
  \node[draw=zolive,rounded corners=2pt,minimum width=14mm,minimum height=7mm,fill=zsoft] (c) at (70,0) {queda};
  \draw[-{Latex[length=2mm]},zmute] (a) -- node[above,font=\scriptsize]{vende $\tfrac12$ + 2} (b);
  \draw[-{Latex[length=2mm]},zmute] (b) -- node[above,font=\scriptsize]{…} (c);
  \draw[-{Latex[length=2mm]},zolive] (c) to[bend left=25] node[below,font=\scriptsize]{al revés: $+2$ y luego $\times 2$} (a);
\end{tikzpicture}
\end{center}

\subsubsection{La barra que se parte}
Cuando los datos son fracciones de lo que queda, una \clave{barra} ordena todo. Si alguien gasta $\tfrac13$ y luego
$\tfrac14$ del resto, se dibuja la barra en tercios, se tacha uno, y lo que queda (dos tercios) se parte en cuartos.
Como $\tfrac23$ partido en 4 da pedazos de $\tfrac16$, conviene partir toda la barra en sextos desde el comienzo.

\begin{center}
\begin{tikzpicture}[x=9mm,y=7mm,font=\scriptsize]
  \foreach \i in {0,...,5}{\draw[zink,line width=.4pt,fill=\ifnum\i<2 zcoral!25\else\ifnum\i<3 zsky!30\else white\fi\fi] (\i,0) rectangle (\i+1,1);}
  \draw[zink,line width=1.1pt] (0,0) rectangle (6,1); \draw[zink,line width=1.1pt] (2,0)--(2,1);
  \draw[decorate,decoration={brace,amplitude=4pt}] (0,1.1) -- node[above=3pt]{$\tfrac13$ gastado} (2,1.1);
  \draw[decorate,decoration={brace,amplitude=4pt}] (2,1.1) -- node[above=3pt]{$\tfrac14$ del resto} (3,1.1);
  \draw[decorate,decoration={brace,amplitude=4pt,mirror}] (3,-.1) -- node[below=3pt]{queda: $\tfrac36 = \tfrac12$} (6,-.1);
\end{tikzpicture}
\end{center}

\begin{regla}
\begin{itemize}
  \item Para deshacer un paso, invierte la operación y el orden: si el paso fue «dividir entre 2 y restar 2», al revés es
  «sumar 2 y multiplicar por 2».
  \item Si se quita $\tfrac ab$ de lo que hay, queda $\tfrac{b-a}{b}$; si después se quita $\tfrac cd$ del resto, queda
  $\tfrac{b-a}{b}\cdot\tfrac{d-c}{d}$ del inicio.
  \item Cuando se conoce una parte y su fracción, la unidad sale dividiendo: si $\tfrac{3}{8}$ del total son 45, el total
  es $45 \div \tfrac38 = 120$.
\end{itemize}
\end{regla}

\subsection{Ejemplos resueltos}
\begin{ejemplo}[Tres paradas]
Una persona gasta en tres tiendas, en cada una la mitad de lo que lleva más \$1\,000. Sale de la tercera sin dinero.
¿Cuánto llevaba al empezar?
\solucion De atrás hacia adelante. Antes de la tercera tienda: $(0 + 1\,000)\cdot 2 = 2\,000$. Antes de la segunda:
$(2\,000 + 1\,000)\cdot 2 = 6\,000$. Antes de la primera: $(6\,000 + 1\,000)\cdot 2 = 14\,000$.
\paso{Comprobación.} $14\,000 \to 14\,000 - 8\,000 = 6\,000 \to 6\,000 - 4\,000 = 2\,000 \to 2\,000 - 2\,000 = 0$.
\end{ejemplo}

\begin{ejemplo}[La herencia]
Un padre deja $\tfrac25$ de su herencia al hijo mayor, $\tfrac13$ del resto al segundo y lo demás, \$24 millones, al
menor. ¿De cuánto era la herencia?
\solucion Tras el mayor queda $\tfrac35$; el segundo se lleva $\tfrac13$ de eso, así que el menor recibe
$\tfrac23 \cdot \tfrac35 = \tfrac25$ de la herencia. Si $\tfrac25$ son 24 millones, $\tfrac15$ son 12 y la herencia
es 60 millones. Curiosidad: el mayor y el menor reciben lo mismo.
\end{ejemplo}

\begin{ejemplo}[El rebote]
Una pelota rebota hasta $\tfrac23$ de la altura desde la que cae. Después del tercer rebote sube 32 cm. ¿Desde qué
altura se soltó?
\solucion Cada rebote multiplica por $\tfrac23$; al revés, cada paso multiplica por $\tfrac32$:
$32 \cdot \tfrac32 = 48$, $48 \cdot \tfrac32 = 72$, $72 \cdot \tfrac32 = 108$ cm.
\end{ejemplo}

\begin{trampa}
Al deshacer, el orden importa. Si el paso fue «la mitad \emph{más} 2», lo que queda es «la mitad \emph{menos} 2» del
número. Quien deshace multiplicando primero por 2 y luego sumando 2 obtiene otra respuesta (que suele ser una opción).
\end{trampa}

\begin{atajo}
Si las opciones son números concretos, haz la historia hacia adelante con cada opción: en menos de un minuto sabes cuál
termina en el dato final.
\end{atajo}

\taller
\piso{Calienta}{hacia atrás}
\begin{ejercicios}
\item Pienso un número, lo multiplico por 3, le resto 7 y obtengo 32. ¿Qué número pensé? \sol{$(32 + 7) \div 3 = 13$.}
\item Gasté $\tfrac34$ de mi dinero y me quedan \$6\,000. ¿Cuánto tenía? \sol{$\tfrac14$ son 6\,000: tenía 24\,000.}
\item De un tanque se saca $\tfrac15$ y luego $\tfrac14$ del resto. ¿Qué fracción queda? \sol{$\tfrac45 \cdot \tfrac34 = \tfrac35$.}
\end{ejercicios}
\piso{Entrena}{problemas}
\begin{ejercicios}
\item Un ciclista recorre el primer día $\tfrac13$ de una ruta, el segundo $\tfrac12$ de lo que falta y el tercero los 40 km restantes. ¿Cuánto mide la ruta? \sol{Queda $\tfrac23 \cdot \tfrac12 = \tfrac13$; la ruta mide 120 km.}
\item Laura regala la mitad de sus láminas y 3 más; luego la mitad de las que le quedan y 3 más, y le quedan 6. ¿Cuántas tenía? \sol{Atrás: $(6+3)\cdot2 = 18$; $(18+3)\cdot 2 = 42$.}
\item En una fiesta se comen $\tfrac38$ de una torta y luego $\tfrac35$ de lo que sobró. ¿Qué fracción de la torta quedó? \sol{$\tfrac58 \cdot \tfrac25 = \tfrac14$.}
\end{ejercicios}
\piso{Como en el examen}{opción múltiple}
\begin{ejercicios}
\item Un estudiante gasta $\tfrac{2}{7}$ de su mesada en transporte y $\tfrac{3}{5}$ del resto en comida. Le quedan \$40\,000. ¿Cuál es su mesada?
\opciones{\$100\,000}{\$140\,000}{\$120\,000}{\$175\,000}
\sol{B. Queda $\tfrac57 \cdot \tfrac25 = \tfrac27$; $\tfrac27$ son 40\,000, así que $\tfrac17$ es 20\,000 y la mesada 140\,000.}
\item Un número se duplica, al resultado se le suma 6, se divide entre 4 y se obtiene 9. El número es
\opciones{12}{15}{18}{21}
\sol{B. Atrás: $9 \cdot 4 = 36$, $36 - 6 = 30$, $30 \div 2 = 15$.}
\item Una bacteria se divide en dos cada hora. A las 10 a.\,m. un frasco está lleno. ¿A qué hora estaba lleno hasta la cuarta parte?
\opciones{4:00 a.\,m.}{8:00 a.\,m.}{7:30 a.\,m.}{9:00 a.\,m.}
\sol{B. Hacia atrás, cada hora antes la mitad: a las 9 la mitad, a las 8 la cuarta parte.}
\end{ejercicios}

\begin{retoresuelto}
Hacia atrás: antes de la tarde tenía $(5 + 2)\cdot 2 = 14$; antes de la mañana, $(14 + 2) \cdot 2 = 32$.
Comprobación: vende $16 + 2 = 18$ y le quedan 14; vende $7 + 2 = 9$ y le quedan 5. \textbf{Opción B.}
\end{retoresuelto}
\bitacora

% ------------------------------------------------------------------
\leccion{pu-ar-2}{Promedios}{Juegos de las torres que se nivelan y la balanza del promedio.}

\begin{reto}
El promedio de 5 números es 18. Si se quita uno de ellos, el promedio de los otros cuatro es 16. ¿Qué número se quitó?
\opciones{2}{20}{26}{34}
\end{reto}

\subsection{La idea: el promedio nivela}
Imagina torres de cubos de distintas alturas. Si se pasan cubos de las altas a las bajas hasta que todas queden iguales,
esa altura común es el \clave{promedio}. Nada se pierde ni se crea: el total de cubos es el mismo.

\begin{center}
\begin{tikzpicture}[x=6mm,y=4mm]
  \foreach \x/\h in {0/7,1/3,2/5,3/1,4/4}{\foreach \k in {1,...,\h}{\draw[zink,line width=.3pt,fill=zlime!60] (\x*1.2,\k-1) rectangle ++(1,1);}}
  \draw[zcoral,dashed,line width=.8pt] (-.3,4) -- (6.1,4) node[right,font=\scriptsize,text=zcoral]{promedio $= 20 / 5 = 4$};
\end{tikzpicture}
\end{center}

De ahí la regla más útil del tema: \textbf{total $=$ promedio $\times$ cantidad}. Casi todo problema de promedios se
resuelve pasando a totales, operando con totales y volviendo a promediar.

\subsubsection{Desviaciones que se cancelan}
Otra forma de verlo: las distancias al promedio se compensan. En el dibujo, $7$ sobra 3, $5$ sobra 1, $3$ falta 1, $1$
falta 3 y $4$ está justo: $+3 + 1 - 1 - 3 + 0 = 0$. Esto permite calcular promedios eligiendo una «base» cómoda:
el promedio de 47, 52, 49 y 56 es $50 + \tfrac{-3 + 2 - 1 + 6}{4} = 50 + 1 = 51$.

\subsubsection{Promedio ponderado}
Si los grupos tienen tamaños distintos, no se promedian los promedios: se pesan. Un curso de 20 con promedio 3.5 y otro
de 30 con promedio 4.0 tienen un promedio conjunto de $\tfrac{20 \cdot 3.5 + 30 \cdot 4.0}{50} = \tfrac{70 + 120}{50}
= 3.8$, más cerca de 4.0 porque ese grupo pesa más.

\begin{regla}
\[\bar x = \frac{\text{suma}}{\text{cantidad}}\qquad\text{suma} = \bar x \cdot n\qquad
\bar x_{\text{conjunto}} = \frac{n_1\bar x_1 + n_2 \bar x_2}{n_1 + n_2}\]
\begin{itemize}
  \item Si a cada dato se le suma $k$, el promedio aumenta $k$; si cada dato se multiplica por $k$, el promedio también.
  \item El promedio de números consecutivos (o en progresión aritmética) es el del centro, o el punto medio entre el
  primero y el último.
\end{itemize}
\end{regla}

\subsection{Ejemplos resueltos}
\begin{ejemplo}[La nota que falta]
Las notas de Sofía en cuatro quices son 3.2, 4.0, 3.6 y 4.4. ¿Cuánto necesita en el quinto para que el promedio sea 4.0?
\solucion Con promedio 4.0 en cinco quices, la suma debe ser $20$. Ya tiene $3.2 + 4.0 + 3.6 + 4.4 = 15.2$. Le faltan
$20 - 15.2 = 4.8$.
\end{ejemplo}

\begin{ejemplo}[Grupos que se juntan]
El peso promedio de 12 hombres es 70 kg y el de 8 mujeres es 60 kg. ¿Cuál es el promedio del grupo?
\solucion Totales: $12 \cdot 70 = 840$ y $8 \cdot 60 = 480$. Promedio: $\tfrac{1\,320}{20} = 66$ kg. No es 65.
\end{ejemplo}

\begin{ejemplo}[Consecutivos]
El promedio de siete números enteros consecutivos es 31. ¿Cuál es el mayor?
\solucion Con siete consecutivos, el promedio es el del centro: 31 es el cuarto. Son 28, 29, 30, 31, 32, 33, 34. El
mayor es 34.
\end{ejemplo}

\begin{ejemplo}[Un dato mal copiado]
El promedio de 10 datos es 24, pero uno se copió como 42 cuando era 12. ¿Cuál es el promedio correcto?
\solucion La suma bajó en $42 - 12 = 30$, así que el promedio baja $\tfrac{30}{10} = 3$: el correcto es 21.
\end{ejemplo}

\begin{trampa}
Promediar promedios sin pesar. «El promedio de 3.5 y 4.0 es 3.75» solo es cierto si los dos grupos son del mismo tamaño.
\end{trampa}

\begin{atajo}
Para promediar números grandes, elige una base cercana, promedia solo las diferencias y súmalas a la base.
\end{atajo}

\taller
\piso{Calienta}{totales}
\begin{ejercicios}
\item Promedio de 12, 15, 18, 21 y 24. \sol{18 (el del centro).}
\item Si el promedio de 6 números es 9, ¿cuánto suman? \sol{54.}
\item Promedio de 98, 103, 101 y 106 (usa base 100). \sol{$100 + \tfrac{-2+3+1+6}{4} = 102$.}
\end{ejercicios}
\piso{Entrena}{problemas}
\begin{ejercicios}
\item El promedio de edad de 4 hermanos es 15. Si se incluye la madre, el promedio sube a 20. ¿Cuántos años tiene la madre? \sol{$5 \cdot 20 - 4 \cdot 15 = 100 - 60 = 40$.}
\item Un curso de 25 estudiantes sacó 3.2 de promedio y otro de 15 sacó 4.0. ¿Promedio de los 40? \sol{$\tfrac{80 + 60}{40} = 3.5$.}
\item Un bus recorre 120 km a 60 km/h y vuelve a 40 km/h. ¿Cuál es su rapidez promedio en todo el viaje? \sol{Tiempo: $2 + 3 = 5$ h para 240 km: 48 km/h (no 50).}
\end{ejercicios}
\piso{Como en el examen}{opción múltiple}
\begin{ejercicios}
\item El promedio de $a$, $b$ y $c$ es 10 y el de $a$ y $b$ es 7. ¿Cuánto vale $c$?
\opciones{3}{13}{16}{17}
\sol{C. $3 \cdot 10 - 2 \cdot 7 = 30 - 14 = 16$.}
\item Si a cada uno de 8 datos se le suma 5, el nuevo promedio
\opcionesd{aumenta 5}{aumenta 40}{aumenta $\tfrac58$}{no cambia}
\sol{A. Cada dato sube 5; la suma sube 40 y el promedio $\tfrac{40}{8} = 5$.}
\item En un salón de 30 estudiantes el promedio fue 3.6. Los 12 hombres promediaron 3.0. ¿Cuál fue el promedio de las mujeres?
\opciones{4.0}{4.2}{3.9}{4.4}
\sol{A. Total 108, hombres 36, mujeres $72 / 18 = 4.0$.}
\end{ejercicios}

\begin{retoresuelto}
Total de los cinco: $5 \cdot 18 = 90$. Total de los cuatro: $4 \cdot 16 = 64$. El que se quitó es $90 - 64 = 26$.
\textbf{Opción C.}
\end{retoresuelto}
\bitacora

% ------------------------------------------------------------------
\leccion{pu-ar-3}{Criptoaritmética}{Juego de las cajas con cifras escondidas.}

\begin{reto}
En la suma
\[\begin{array}{r}\mathrm{S\,E\,N\,D}\\+\;\mathrm{M\,O\,R\,E}\\\hline\mathrm{M\,O\,N\,E\,Y}\end{array}\]
cada letra es una cifra distinta y ninguna palabra empieza por cero. ¿Cuánto vale $M$?
\opciones{0}{1}{2}{9}
\end{reto}

\subsection{La idea: las columnas hablan}
Una cuenta con letras o con huecos se resuelve \clave{columna por columna}, y casi siempre se empieza por los extremos:
la cifra de la izquierda (que no puede ser 0) y las llevadas. Tres hechos sostienen todo:
\begin{enumerate}
  \item Al sumar dos números, lo que se lleva a la columna siguiente es 0 o 1 (nunca más, porque $9 + 9 + 1 = 19$).
  \item Si el resultado tiene una cifra más que los sumandos, esa cifra nueva es 1.
  \item En una multiplicación, la última cifra del producto depende solo de las últimas cifras de los factores.
\end{enumerate}

\begin{regla}[Estrategia]
\begin{enumerate}
  \item Busca la columna con menos incógnitas o con una llevada obligada.
  \item Anota las llevadas encima, como en la escuela.
  \item Cuando quedan pocas opciones, \emph{prueba} y descarta. Escribe lo que vas descartando.
\end{enumerate}
\end{regla}

\subsection{Ejemplos resueltos}
\begin{ejemplo}[Cifras escondidas]
Completa: $\;\square\,4\,\square + 3\,\square\,8 = 1\,0\,0\,3$.
\solucion \paso{Unidades.} $\square + 8$ termina en 3: el hueco es 5 y se lleva 1.
\paso{Decenas.} $4 + \square + 1$ termina en 0: el hueco es 5 y se lleva 1.
\paso{Centenas.} $\square + 3 + 1 = 10$: el hueco es 6.
La cuenta es $645 + 358 = 1\,003$. Cada columna dejó una sola posibilidad porque se empezó por la derecha, donde no hay
llevadas desconocidas.
\end{ejemplo}

\begin{ejemplo}[Una multiplicación]
Halla el número $\overline{ABCD}$ de cifras distintas tal que $\overline{ABCD} \cdot 4 = \overline{DCBA}$.
\solucion
\paso{La primera cifra.} El producto sigue teniendo cuatro cifras, así que $A \le 2$. Además $A$ es la última cifra de
$4D$, que es par: $A = 2$.
\paso{La última.} $4D$ termina en 2: $D = 3$ u $8$. Como $D \ge 4A = 8$, $D = 8$. Las unidades dan $32$: se lleva 3.
\paso{Los millares.} $4 \cdot 2 = 8 = D$, así que de las centenas no se lleva nada: $4B$ más lo que se lleve es menor
que 10, y $B \le 2$.
\paso{Las decenas.} $4C + 3$ es impar y termina en $B$: $B$ es impar, así que $B = 1$. Entonces $4C$ termina en 8:
$C = 2$ o $C = 7$; $C = 2$ repetiría la cifra de $A$, así que $C = 7$. Comprobación en las centenas: $4 \cdot 1 + 3 = 7 = C$.
El número es $2\,178$, y en efecto $2\,178 \cdot 4 = 8\,712$.
\end{ejemplo}

\begin{ejemplo}[Suma de dos cifras repetidas]
Si $\overline{AA} + \overline{BB} = \overline{1C0}$ con $A$, $B$, $C$ cifras distintas y distintas de cero, ¿cuánto
vale $A + B$?
\solucion $\overline{AA} = 11A$ y $\overline{BB} = 11B$, así que la suma es $11(A + B)$, un múltiplo de 11 entre 100
y 199 que termina en 0: solo $110$. Entonces $A + B = 10$ y $C = 1$.
\end{ejemplo}

\begin{trampa}
Olvidar que la cifra inicial de un número no puede ser 0, o que las letras distintas son cifras distintas. Esas dos
condiciones son las que dejan una sola solución.
\end{trampa}

\begin{atajo}
Escribir el número por posiciones: $\overline{abc} = 100a + 10b + c$. Muchos acertijos de cifras se vuelven una ecuación
corta.
\end{atajo}

\taller
\piso{Calienta}{huecos}
\begin{ejercicios}
\item Completa $3\square 6 + 2 5\square = 6 0 3$. \sol{$346 + 257 = 603$: los huecos son 4 y 7.}
\item Completa $5\square 3 - \square 7 5 = 2 6 8$. \sol{$543 - 275 = 268$: los huecos son 4 y 2.}
\item Si $\overline{ab} + \overline{ba} = 121$, ¿cuánto vale $a + b$? \sol{$11(a + b) = 121$: $a + b = 11$.}
\end{ejercicios}
\piso{Entrena}{problemas}
\begin{ejercicios}
\item Si $\overline{ab} - \overline{ba} = 45$, ¿cuánto vale $a - b$? \sol{$9(a - b) = 45$: $a - b = 5$.}
\item En $\overline{A\,B} + \overline{A\,B} + \overline{A\,B} = \overline{C\,B\,B}$ con cifras distintas, halla $\overline{AB}$. \sol{$3B$ termina en $B$: $B = 0$ o $5$. Con $B=5$: $3 \cdot \overline{A5}$ debe terminar en 55: $\overline{A5} = 85$ da 255. $\overline{AB} = 85$, $C = 2$.}
\item ¿Cuántos números de dos cifras aumentan en 27 al invertir sus cifras? \sol{$\overline{ba} - \overline{ab} = 9(b - a) = 27$: $b = a + 3$, con $a \ge 1$: 14, 25, 36, 47, 58, 69. Son 6.}
\end{ejercicios}
\piso{Como en el examen}{opción múltiple}
\begin{ejercicios}
\item Si $A$ es una cifra par y $\overline{2A} \cdot A = \overline{1BA}$, ¿cuánto es $A + B$?
\opciones{9}{11}{13}{6}
\sol{B. $A \cdot A$ termina en $A$ y $A$ es par: $A = 6$. $26 \cdot 6 = 156$, así que $B = 5$ y $A + B = 11$.}
\item Se borra una cifra de un número de 4 cifras: $8\,\square\,3\,6$. Si el número es múltiplo de 9, la cifra borrada es
\opciones{1}{3}{0}{5}
\sol{A. $8 + 3 + 6 = 17$; para llegar a 18 falta 1.}
\item Al sumar $\overline{ABC} + \overline{CBA}$ con $A > C$ se obtiene 888. ¿Cuántos valores puede tomar $B$?
\opciones{1}{2}{4}{5}
\sol{A. Las unidades dan $A + C = 8$ (no puede ser 18, porque $A \ne C$). En las decenas, $2B$ termina en 8: $B = 4$ o $B = 9$; con $B = 9$ se llevaría 1 y las centenas darían $A + C = 7$, contradicción. Solo $B = 4$.}
\end{ejercicios}

\begin{retoresuelto}
El resultado MONEY tiene una cifra más que los sumandos, así que $M$ es lo que se lleva de la última columna: $M = 1$.
(Esta suma, famosa desde 1924, tiene solución única: $9567 + 1085 = 10652$.) \textbf{Opción B.}
\end{retoresuelto}
\bitacora

% ------------------------------------------------------------------
\leccion{pu-ar-4}{Máximos y mínimos}{Juegos del cajón de medias y del peor caso.}

\begin{reto}
En un cajón oscuro hay 6 medias rojas, 8 azules y 10 negras. ¿Cuántas hay que sacar, como mínimo, para estar
\emph{seguros} de tener dos del mismo color?
\opciones{2}{3}{4}{11}
\end{reto}

\subsection{La idea: pensar en el peor caso}
La palabra «\emph{seguro}» cambia el problema: no importa qué puede pasar con suerte, sino qué pasa con la \clave{peor
suerte posible}. Para asegurar algo, imagina a un adversario que saca las medias en el orden que más te perjudica y
cuenta cuántas necesita para no poder evitar lo que pides.

Este razonamiento se llama \clave{principio del palomar}: si hay más palomas que casillas, alguna casilla tiene al menos
dos palomas. Con $n$ casillas, $n + 1$ objetos garantizan una casilla repetida; $kn + 1$ objetos garantizan una casilla
con $k + 1$.

\subsubsection{Máximos y mínimos con condiciones}
La otra cara del tema: cuando se reparte un total con condiciones («todos distintos», «cada uno al menos 2»), para
\emph{maximizar} a uno se \emph{minimiza} a los demás, y al revés.

\begin{regla}
\begin{itemize}
  \item Para garantizar $k$ objetos de un mismo tipo entre $n$ tipos: en el peor caso se sacan $k - 1$ de cada tipo; uno
  más asegura el resultado. Mínimo: $n(k - 1) + 1$ (si hay suficientes de cada tipo).
  \item Para garantizar $k$ de un tipo \emph{específico}: en el peor caso salen primero \emph{todos} los de los otros tipos.
  \item Para hacer grande una parte de un total fijo, haz lo más pequeñas posibles las otras.
\end{itemize}
\end{regla}

\subsection{Ejemplos resueltos}
\begin{ejemplo}[Un color específico]
Con el cajón del reto (6 rojas, 8 azules, 10 negras), ¿cuántas hay que sacar para asegurar dos rojas?
\solucion Peor caso: salen las 8 azules, las 10 negras y una roja: 19 medias sin tener dos rojas. La siguiente es roja
seguro: hacen falta $20$.
\end{ejemplo}

\begin{ejemplo}[Cumpleaños]
¿Cuántas personas debe haber en un grupo para asegurar que al menos tres cumplen años el mismo mes?
\solucion Casillas: 12 meses. Para que ningún mes tenga tres, a lo sumo 2 por mes: 24 personas. Con $25$ es seguro.
\end{ejemplo}

\begin{ejemplo}[Repartir para maximizar]
Cinco amigos tienen en total 30 láminas, todos cantidades distintas y cada uno al menos una. ¿Cuántas puede tener, como
máximo, el que más tiene?
\solucion Los otros cuatro, lo menos posible y distintos: $1 + 2 + 3 + 4 = 10$. Al mayor le quedan $30 - 10 = 20$.
\end{ejemplo}

\begin{ejemplo}[Lo mínimo garantizado]
En un curso de 40 estudiantes, 28 aprobaron matemáticas y 25 aprobaron lenguaje. ¿Cuántos aprobaron las dos, como mínimo?
\solucion $28 + 25 = 53$, pero solo hay 40 personas: al menos $53 - 40 = 13$ se contaron dos veces. Mínimo 13 (y como
máximo 25).
\end{ejemplo}

\begin{trampa}
Responder con el caso de buena suerte (2 medias) cuando la pregunta dice «para estar seguros». Esa opción casi siempre
está.
\end{trampa}

\begin{atajo}
Escribe una línea con el peor caso («salen 1 roja, 1 azul, 1 negra…») y suma 1. Si la condición es sobre un tipo
específico, el peor caso agota los demás tipos primero.
\end{atajo}

\taller
\piso{Calienta}{peor caso}
\begin{ejercicios}
\item En una bolsa hay 5 bolas blancas y 7 negras. ¿Cuántas sacar para asegurar una negra? \sol{$5 + 1 = 6$.}
\item ¿Cuántas cartas de una baraja de 4 palos hay que sacar para asegurar 3 del mismo palo? \sol{$4 \cdot 2 + 1 = 9$.}
\item Tres números naturales distintos suman 20. ¿Cuál es el mayor valor posible del más grande? \sol{$20 - 1 - 2 = 17$.}
\end{ejercicios}
\piso{Entrena}{problemas}
\begin{ejercicios}
\item En una caja hay 4 pares de guantes negros y 3 pares cafés, revueltos. ¿Cuántos guantes sacar para asegurar un par que sirva (una mano izquierda y una derecha del mismo color)? \sol{Peor caso: salen las 7 manos izquierdas (o las 7 derechas). El siguiente completa un par: 8.}
\item En un salón de 35 estudiantes, ¿se puede asegurar que 3 nacieron el mismo día de la semana? \sol{Sí: $7 \cdot 2 = 14 < 35$. De hecho se aseguran 5: $7\cdot 4 = 28 < 35$.}
\item Cuatro números enteros positivos tienen promedio 10. ¿Cuál es el mayor valor que puede tomar uno de ellos? \sol{Suma 40; los otros tres valen 1: el mayor puede ser 37.}
\end{ejercicios}
\piso{Como en el examen}{opción múltiple}
\begin{ejercicios}
\item Una urna tiene 10 bolas rojas, 10 verdes y 10 amarillas. ¿Cuántas sacar para asegurar 4 del mismo color?
\opciones{4}{10}{12}{13}
\sol{B. Peor caso: 3 de cada color (9 bolas). La décima asegura 4 de un color.}
\item De 50 personas, 32 hablan inglés y 27 francés. El número mínimo de personas que hablan ambos idiomas es
\opciones{5}{9}{18}{27}
\sol{B. $32 + 27 - 50 = 9$.}
\item Seis números enteros positivos distintos suman 40. El menor valor posible del mayor es
\opciones{9}{10}{8}{35}
\sol{B. Para que el mayor sea pequeño, los números deben estar cerca entre sí. Con mayor 9, el máximo posible es $4+5+6+7+8+9 = 39 < 40$. Con mayor 10: $4+5+6+7+8+10 = 40$. Respuesta 10.}
\end{ejercicios}

\begin{retoresuelto}
Peor caso: una roja, una azul y una negra (tres medias sin repetir). La cuarta repite color sin remedio.
\textbf{Opción C.} La A es la respuesta con suerte, y la D confunde con «asegurar una de cada color».
\end{retoresuelto}
\bitacora

% ------------------------------------------------------------------
\evaluacion
\begin{ejercicios}
\item María gasta $\tfrac{1}{4}$ de su sueldo en arriendo y $\tfrac{1}{3}$ del resto en mercado. Si le quedan \$900\,000, su sueldo es
\opciones{\$1\,500\,000}{\$1\,800\,000}{\$2\,000\,000}{\$1\,600\,000}
\sol{B. Queda $\tfrac34 \cdot \tfrac23 = \tfrac12$. (2.1)}
\item Pienso un número, le sumo 5, lo multiplico por 3 y le resto 6. Obtengo 33. El número es
\opciones{8}{6}{10}{13}
\sol{A. $(33 + 6)\div 3 - 5 = 8$. (2.1)}
\item Una pelota rebota hasta la mitad de la altura de la que cae. Si después del cuarto rebote sube 5 cm, se soltó desde
\opciones{40 cm}{80 cm}{100 cm}{160 cm}
\sol{B. $5 \cdot 2^4 = 80$. (2.1)}
\item El promedio de 6 números es 12. Si se agrega el 26, el nuevo promedio es
\opciones{14}{13}{19}{15}
\sol{A. $(72 + 26)/7 = 14$. (2.2)}
\item Diego tiene promedio 3.8 en 4 notas. ¿Qué nota necesita en la quinta para subir a 4.0?
\opciones{4.2}{4.6}{4.8}{5.0}
\sol{C. $20 - 15.2 = 4.8$. (2.2)}
\item En un grupo, 10 personas promedian 1.70 m y otras 30 promedian 1.60 m. El promedio de todo el grupo es
\opciones{1.65 m}{1.625 m}{1.675 m}{1.63 m}
\sol{B. $(17 + 48)/40 = 1.625$. (2.2)}
\item Si $\overline{ab} + \overline{ba} = 154$, entonces $a + b$ vale
\opciones{11}{14}{13}{15}
\sol{B. $11(a + b) = 154$. (2.3)}
\item Si $\overline{3A5}$ es múltiplo de 9, $A$ vale
\opciones{1}{4}{0}{9}
\sol{A. $3 + A + 5$ debe ser 9 o 18: $A = 1$ o $A = 10$ (no es cifra). (2.3)}
\item Un número de dos cifras disminuye en 36 al invertir sus cifras, y sus cifras suman 12. El número es
\opciones{84}{48}{73}{95}
\sol{A. $9(a - b) = 36$ y $a + b = 12$: $a = 8$, $b = 4$. (2.3)}
\item Una caja tiene 7 lápices rojos, 5 azules y 3 verdes. ¿Cuántos sacar a ciegas para asegurar uno de cada color?
\opciones{3}{8}{13}{15}
\sol{C. Peor caso: salen los 7 rojos y los 5 azules (12); el 13.º es verde. (2.4)}
\item En un club de 60 socios, 45 juegan tenis y 38 natación. Como mínimo, ¿cuántos practican los dos deportes?
\opciones{15}{23}{38}{22}
\sol{B. $45 + 38 - 60 = 23$. (2.4)}
\end{ejercicios}
